\documentclass{article}
\usepackage{fleqn}
\usepackage{epsf}
\usepackage{aima2e-slides}

\begin{document}

\begin{huge}
\titleslide{Học từ các quan sát (Learning from Observations)}{Chương 18, Phần 1--3}

\sf

%%%%%%%%%%%% Slide %%%%%%%%%%%%%%%%%%%%%%%%%%%%%%%%%%%%%%%%%%%%%%%%%%%
\heading{Phác thảo}

\blob Tác nhân học tập

\blob Học quy nạp

\blob Học cây quyết định

\blob Đo lường hiệu quả học tập


%%%%%%%%%%%% Slide %%%%%%%%%%%%%%%%%%%%%%%%%%%%%%%%%%%%%%%%%%%%%%%%%%%
\heading{Học}

Học tập là điều cần thiết cho những môi trường chưa biết,\\
tức là khi nhà thiết kế thiếu sự toàn tri

Học tập hữu ích như một phương pháp xây dựng hệ thống,\\
tức là, cho tác nhân tiếp xúc với thực tế thay vì cố gắng viết nó ra

Việc học sửa đổi cơ chế quyết định của tác nhân để cải thiện hiệu suất


%%%%%%%%%%%% Slide %%%%%%%%%%%%%%%%%%%%%%%%%%%%%%%%%%%%%%%%%%%%%%%%%%%
\heading{Tác nhân học tập}

\epsfxsize=0,95\textwidth
\fig{\file{figures}{learning-model.ps}}

% Points to make:
% 1) Performance element is what we have called ``agent'' up to now
% 2) Critic/LearningElement/ProblemGenerator do the ``improving''
% 3) Performance standard is fixed; 
%    a) can't adjust performance standard to flatter own behaviour
%    b) no standard *in the environment*: ordinary chess and
%       suicide chess LOOK identical. Essentially, certain kinds of percepts
%       are ``hardwired'' as good/bad (e.g., pain, hunger)
% 4) Learning element may use knowledge already acquired in the perf. element
% 5) Learning may require experimentation - actions an agent might not
%    normally consider such as dropping rocks for the Tower of Pisa



%%%%%%%%%%%% Slide %%%%%%%%%%%%%%%%%%%%%%%%%%%%%%%%%%%%%%%%%%%%%%%%%%%
\heading{Phần tử học tập}

Thiết kế thành phần học tập được quyết định bởi\al
\blob loại phần tử hiệu suất nào được sử dụng\al
\blob thành phần chức năng nào sẽ được học\al
\blob cách thể hiện thành phần chức năng đó\al
\blob có những loại phản hồi nào \\
Các tình huống ví dụ:

\vspace*{0.1in}

\epsfxsize=1.07\textwidth
\fig{\file{figures}{learning-elements.ps}}

\defn{Học có giám sát}: câu trả lời đúng cho từng trường hợp\\
\defn{Học tăng cường}: phần thưởng không thường xuyên

% Points to make:
% 1) Learning transition model is ``supervised'' if observable
% 2) Supervised learning of correct actions requires ``teacher''
% 3) Reinforcement learning is harder, but requires no teacher


%%%%%%%%%%%% Slide %%%%%%%%%%%%%%%%%%%%%%%%%%%%%%%%%%%%%%%%%%%%%%%%%%%
\heading{Học quy nạp (còn gọi là~Khoa học)}

Hình thức đơn giản nhất: học hàm từ các ví dụ (\emph{tabula rasa})

\al là \emph{hàm đích}

Một \defn{example} là một cặp \mat{$x$}, \mat{$f$}, ví dụ: 
\mat{$\begin{array}{c|c|c}O&O&X \\
                     \hline
                      &X&  \\
                     \hline
                     X& &  \end{array}\ \ ,\ \ +1$}

Bài toán: tìm một(n) \defn{giả thuyết} \mat{$h$}\nl
   sao cho \mat{$h \approx f$}\nl
   đưa ra một \defn{tập huấn luyện} gồm các ví dụ

(\emph{Đây là mô hình học tập thực tế được đơn giản hóa cao:\al
  -- Bỏ qua kiến thức trước đây\al
  -- Giả sử một ``môi trường'' xác định, có thể quan sát được \al
  -- Giả sử các ví dụ được cho là \emph{}\al
  -- Giả sử rằng tác nhân \emph{muốn} tìm hiểu \mat{$f$}---tại sao?})

% Points to make:
% 1) For Science, the ``agent'' is the whole human race over all time
% 2) tabula rasa = ``blank slate'' (lit. erased writing tablet)



%%%%%%%%%%%% Slide %%%%%%%%%%%%%%%%%%%%%%%%%%%%%%%%%%%%%%%%%%%%%%%%%%%
\heading{Phương pháp học quy nạp}

Xây dựng/điều chỉnh \mat{$h$} để phù hợp với \mat{$f$} trên tập huấn luyện \\
(\mat{$h$} là \defn{nhất quán} nếu nó đồng ý với \mat{$f$} trên tất cả các ví dụ)

Ví dụ: khớp đường cong:

\vspace*{0.2in}

\epsfxsize=0,5\textwidth
\fig{\sfile{figures}{curve-fitting1.ps}}

%%%%%%%%%%%% Slide %%%%%%%%%%%%%%%%%%%%%%%%%%%%%%%%%%%%%%%%%%%%%%%%%%%
\heading{Phương pháp học quy nạp}

Xây dựng/điều chỉnh \mat{$h$} để phù hợp với \mat{$f$} trên tập huấn luyện \\
(\mat{$h$} là \defn{nhất quán} nếu nó đồng ý với \mat{$f$} trên tất cả các ví dụ)

Ví dụ: khớp đường cong:

\vspace*{0.2in}

\epsfxsize=0,5\textwidth
\fig{\sfile{figures}{curve-fitting2.ps}}

%%%%%%%%%%%% Slide %%%%%%%%%%%%%%%%%%%%%%%%%%%%%%%%%%%%%%%%%%%%%%%%%%%
\heading{Phương pháp học quy nạp}

Xây dựng/điều chỉnh \mat{$h$} để phù hợp với \mat{$f$} trên tập huấn luyện \\
(\mat{$h$} là \defn{nhất quán} nếu nó đồng ý với \mat{$f$} trên tất cả các ví dụ)

Ví dụ: khớp đường cong:

\vspace*{0.2in}

\epsfxsize=0,5\textwidth
\fig{\sfile{figures}{curve-fitting3.ps}}

%%%%%%%%%%%% Slide %%%%%%%%%%%%%%%%%%%%%%%%%%%%%%%%%%%%%%%%%%%%%%%%%%%
\heading{Phương pháp học quy nạp}

Xây dựng/điều chỉnh \mat{$h$} để phù hợp với \mat{$f$} trên tập huấn luyện \\
(\mat{$h$} là \defn{nhất quán} nếu nó đồng ý với \mat{$f$} trên tất cả các ví dụ)

Ví dụ: khớp đường cong:

\vspace*{0.2in}

\epsfxsize=0,5\textwidth
\fig{\sfile{figures}{curve-fitting4.ps}}

%%%%%%%%%%%% Slide %%%%%%%%%%%%%%%%%%%%%%%%%%%%%%%%%%%%%%%%%%%%%%%%%%%
\heading{Phương pháp học quy nạp}

Xây dựng/điều chỉnh \mat{$h$} để phù hợp với \mat{$f$} trên tập huấn luyện \\
(\mat{$h$} là \defn{nhất quán} nếu nó đồng ý với \mat{$f$} trên tất cả các ví dụ)

Ví dụ: khớp đường cong:

\vspace*{0.2in}

\epsfxsize=0,5\textwidth
\fig{\sfile{figures}{curve-fitting5.ps}}

%%%%%%%%%%%% Slide %%%%%%%%%%%%%%%%%%%%%%%%%%%%%%%%%%%%%%%%%%%%%%%%%%%
\heading{Phương pháp học quy nạp}

Xây dựng/điều chỉnh \mat{$h$} để phù hợp với \mat{$f$} trên tập huấn luyện \\
(\mat{$h$} là \defn{nhất quán} nếu nó đồng ý với \mat{$f$} trên tất cả các ví dụ)

Ví dụ: khớp đường cong:

\vspace*{0.2in}

\epsfxsize=0,5\textwidth
\fig{\sfile{figures}{curve-fitting5.ps}}

\defn{Dao cạo Ockham}: tối đa hóa sự kết hợp giữa tính nhất quán và sự đơn giản




%%%%%%%%%%%% Slide %%%%%%%%%%%%%%%%%%%%%%%%%%%%%%%%%%%%%%%%%%%%%%%%%%%
\heading{Biểu diễn dựa trên thuộc tính}

Ví dụ được mô tả bởi các giá trị thuộc tính \defn{} (Boolean, rời rạc, liên tục, v.v.)\\
Ví dụ: các tình huống mà tôi sẽ/không đợi bàn:

\input{\file{tables}{restaurant-example-table}}%

\defn{Phân loại} của các ví dụ là \defn{dương} (T) hoặc \defn{âm tính} (F)

% Points to make:
% 1) Give detailed attribute definitions (on p.532)


%%%%%%%%%%%% Slide %%%%%%%%%%%%%%%%%%%%%%%%%%%%%%%%%%%%%%%%%%%%%%%%%%%
\heading{Cây quyết định}

Một đại diện khả dĩ cho các giả thuyết\\
Ví dụ: đây là cây ``true'' để quyết định có nên đợi hay không:

\vspace*{0.2in}

\epsfxsize=0,8\textwidth
\fig{\file{figures}{restaurant-tree.ps}}

% Points to make:
% 1) Explain how the decision tree evaluates an example, e.g., \mat{$X_3$}
% 2) Some of the original set of attributes are irrelevant
% 3) This case is REALIZABLE - i.e., target definable in hypothesis class

%%%%%%%%%%%% Slide %%%%%%%%%%%%%%%%%%%%%%%%%%%%%%%%%%%%%%%%%%%%%%%%%%%
\heading{Tính biểu cảm}

Cây quyết định có thể thể hiện bất kỳ chức năng nào của thuộc tính đầu vào.\\
Ví dụ: đối với các hàm Boolean, đường dẫn hàng bảng chân lý \mat{$\rightarrow$} đến lá:

\vspace*{0.2in}

\epsfxsize=0,65\textwidth
\fig{\file{figures}{xor-decision-tree.ps}}

Điều đáng chú ý là có một cây quyết định nhất quán cho bất kỳ tập huấn luyện nào\\
w/ một đường dẫn đến lá cho mỗi ví dụ (trừ khi \mat{$f$} không xác định trong \mat{$x$})\\
nhưng có lẽ nó sẽ không khái quát hóa được các ví dụ mới

Muốn tìm thêm cây quyết định \emph{compact}


%%%%%%%%%%%% Slide %%%%%%%%%%%%%%%%%%%%%%%%%%%%%%%%%%%%%%%%%%%%%%%%%%%
\heading{Không gian giả thuyết}

\q{Có bao nhiêu cây quyết định riêng biệt với thuộc tính Boolean \mat{$n$}}

% Points to make: get students to think about question before
% going on to next slides which have the answer

%%%%%%%%%%%% Slide %%%%%%%%%%%%%%%%%%%%%%%%%%%%%%%%%%%%%%%%%%%%%%%%%%%
\heading{Không gian giả thuyết}

\q{Có bao nhiêu cây quyết định riêng biệt với thuộc tính Boolean \mat{$n$}}

= số lượng hàm Boolean

%%%%%%%%%%%% Slide %%%%%%%%%%%%%%%%%%%%%%%%%%%%%%%%%%%%%%%%%%%%%%%%%%%
\heading{Không gian giả thuyết}

\q{Có bao nhiêu cây quyết định riêng biệt với thuộc tính Boolean \mat{$n$}}

= số lượng hàm Boolean\\
= số bảng chân lý riêng biệt có hàng \mat{$2^n$}

%%%%%%%%%%%% Slide %%%%%%%%%%%%%%%%%%%%%%%%%%%%%%%%%%%%%%%%%%%%%%%%%%%
\heading{Không gian giả thuyết}

\q{Có bao nhiêu cây quyết định riêng biệt với thuộc tính Boolean \mat{$n$}}

= số lượng hàm Boolean\\
= số bảng chân lý riêng biệt có \mat{$2^n$} hàng = \mat{$2^{2^n}$}

%%%%%%%%%%%% Slide %%%%%%%%%%%%%%%%%%%%%%%%%%%%%%%%%%%%%%%%%%%%%%%%%%%
\heading{Không gian giả thuyết}

\q{Có bao nhiêu cây quyết định riêng biệt với thuộc tính Boolean \mat{$n$}}

= số lượng hàm Boolean\\
= số bảng chân lý riêng biệt có \mat{$2^n$} hàng = \mat{$2^{2^n}$}

Ví dụ: với 6 thuộc tính Boolean, có 18.446.744.073.709.551.616 cây

%%%%%%%%%%%% Slide %%%%%%%%%%%%%%%%%%%%%%%%%%%%%%%%%%%%%%%%%%%%%%%%%%%
\heading{Không gian giả thuyết}

\q{Có bao nhiêu cây quyết định riêng biệt với thuộc tính Boolean \mat{$Hungry\land\lnot Rain$}}

= số lượng hàm Boolean\\
= số bảng chân lý riêng biệt có \mat{$2^n$} hàng = \mat{$2^{2^n}$}

Ví dụ: với 6 thuộc tính Boolean, có 18.446.744.073.709.551.616 cây

\q{Có bao nhiêu giả thuyết liên kết thuần túy (ví dụ: \mat{$Hungry\land\lnot Rain$})}

% Points to make: get students to think about question before
% going on to next slides which have the answer

%%%%%%%%%%%% Slide %%%%%%%%%%%%%%%%%%%%%%%%%%%%%%%%%%%%%%%%%%%%%%%%%%%
\heading{Không gian giả thuyết}

\q{Có bao nhiêu cây quyết định riêng biệt với thuộc tính Boolean \mat{$Hungry\land\lnot Rain$}}

= số lượng hàm Boolean\\
= số bảng chân lý riêng biệt có \mat{$2^n$} hàng = \mat{$2^{2^n}$}

Ví dụ: với 6 thuộc tính Boolean, có 18.446.744.073.709.551.616 cây

\q{Có bao nhiêu giả thuyết liên kết thuần túy (ví dụ: \mat{$Hungry\land\lnot Rain$})}

Mỗi thuộc tính có thể ở (dương), in (âm) hoặc out\nl
\mat{$\implies$} \mat{$3^n$} giả thuyết liên kết khác biệt

Không gian giả thuyết biểu cảm hơn\al
 -- tăng khả năng hàm mục tiêu có thể được biểu thị \quad\smiley\al
 -- tăng số lượng giả thuyết phù hợp với tập huấn luyện\nl
    \mat{$\implies$} có thể nhận được dự đoán tồi tệ hơn\quad\frowny


%%%%%%%%%%%% Slide %%%%%%%%%%%%%%%%%%%%%%%%%%%%%%%%%%%%%%%%%%%%%%%%%%%
\heading{Học cây quyết định}

Mục tiêu: tìm một cây nhỏ phù hợp với các ví dụ đào tạo

Ý tưởng: (đệ quy) chọn thuộc tính ``quan trọng nhất'' làm gốc của cây (phụ)

\code{
\func{Decision-Tree-Learning}{\var{examples{\ac}attributes{\ac}default}}{a decision tree}
    \firstinputs{examples}{set of examples}
    \inputs{attributes}{set of attributes}
    \inputs{default}{default value for the goal predicate}
\bodysep
    \key{if} \var{examples} is empty \key{then} \key{return} \var{default}
    \key{else} \key{if} all \var{examples} have the same classification \key{then} \key{return} the classification
    \key{else} \key{if} \var{attributes} is empty \key{then} \key{return} \prog{Majority-Value}(\var{examples})
    \key{else} 
          \setq{\var{best}}{\prog{Choose-Attribute}(\var{attributes}{\ac}\var{examples})}
          \setq{\var{tree}}{a new decision tree with root test \var{best}}
          \key{for} \key{each} value \var{v}$_i$ of \var{best} \key{do}
                \setq{\var{examples}$_i$}{$\{$elements of \var{examples} with $\var{best} = v_i\}$}
                \setq{subtree}{}\prog{Decision-Tree-Learning}(\var{examples}$_i{\ac}$\var{attributes}$\,-\,$\var{best}{\ac}
                \phantom{\setq{subtree}{}\prog{Decision-Tree-Learning}(}\prog{Majority-Value}(\var{examples}))
                add a branch to \var{tree} with label \var{v}$_i$ and subtree \var{subtree}
          \key{end}
          \key{return} \var{tree}
}


% Points to make:
% 1) Explain inputs
% 2) Explain recursive step -- pick attribute, subdivide examples
% 3) Explain base cases: 
%    a) empty examples - arises for empty branches of non-Boolean parent attribute
%    b) uniform example classification - this is ``normal'' leaf
%    c) empty attributes - target is not deterministic in input attributes

%%%%%%%%%%%% Slide %%%%%%%%%%%%%%%%%%%%%%%%%%%%%%%%%%%%%%%%%%%%%%%%%%%
\heading{Chọn thuộc tính}

Ý tưởng: một thuộc tính tốt sẽ chia các ví dụ thành các tập con
(lý tưởng) ``tất cả đều tích cực'' hoặc ``tất cả tiêu cực''

\vspace*{0.5in}

\epsfxsize=1.0\textwidth
\fig{\file{figures}{restaurant-roots.ps}}

\mat{$Patrons?$} là lựa chọn tốt hơn---cung cấp \emph{thông tin} về phân loại



%%%%%%%%%%%% Slide %%%%%%%%%%%%%%%%%%%%%%%%%%%%%%%%%%%%%%%%%%%%%%%%%%%
\heading{Thông tin}

Thông tin trả lời câu hỏi

Tôi càng không biết gì về câu trả lời ban đầu thì càng
thông tin có trong câu trả lời

Tỷ lệ: 1 bit = câu trả lời cho câu hỏi Boolean có \mat{$\<0.5,0.5\>$} trước

Thông tin trong câu trả lời có trước \mat{$\<P_1,\ldots,P_n\>$} là
\mat{\[
  H(\<P_1,\ldots,P_n\>) = \mysum_{i\eq 1}^n - P_i \log_2 P_i
\]}
(còn gọi là \defn{entropy} trước đó)



%%%%%%%%%%%% Slide %%%%%%%%%%%%%%%%%%%%%%%%%%%%%%%%%%%%%%%%%%%%%%%%%%%
\heading{Thông tin tiếp theo.}

Giả sử chúng ta có \mat{$p$} ví dụ dương và \mat{$n$} âm ở gốc\nl
\mat{$\implies$} \mat{$H(\<p/(p+n),n/(p+n)\>)$} bit cần thiết để phân loại một ví dụ mới \\
Ví dụ: đối với 12 ví dụ về nhà hàng, \mat{$p\eq n\eq 6$} vì vậy chúng tôi cần 1 bit

Một thuộc tính chia các ví dụ \mat{$E$} thành các tập con \mat{$E_i$}, mỗi tập con đó
(chúng tôi hy vọng) cần ít thông tin hơn để hoàn thành việc phân loại

Đặt \mat{$E_i$} có \mat{$p_i$} ví dụ dương và \mat{$n_i$} âm \al
\mat{$\implies$} \mat{$H(\<p_i/(p_i+n_i),n_i/(p_i+n_i)\>)$} bit cần thiết để phân loại một ví dụ mới \al
\mat{$\implies$} \emph{số bit mong đợi} trên mỗi ví dụ trên tất cả các nhánh là
\mat{\[
  \mysum_i \ \ \frac{p_i+n_i}{p+n} \ H(\<p_i/(p_i+n_i),n_i/(p_i+n_i)\>)
\]}
Đối với \mat{$Patrons?$}, đây là 0,459 bit, đối với \mat{$Type$} thì đây là (vẫn) 1 bit

\mat{$\implies$} chọn thuộc tính giảm thiểu thông tin còn lại cần thiết


%%%%%%%%%%%% Slide %%%%%%%%%%%%%%%%%%%%%%%%%%%%%%%%%%%%%%%%%%%%%%%%%%%
\heading{Ví dụ tiếp theo.}

Cây quyết định học được từ 12 ví dụ:

\vspace*{0.2in}

\epsfxsize=0,6\textwidth
\fig{\file{figures}{induced-restaurant-tree.ps}}

Đơn giản hơn nhiều so với cây “đúng”---một giả thuyết phức tạp hơn
không được chứng minh bằng lượng nhỏ dữ liệu



%%%%%%%%%%%% Slide %%%%%%%%%%%%%%%%%%%%%%%%%%%%%%%%%%%%%%%%%%%%%%%%%%%
\heading{Đo hiệu suất}

Làm sao chúng ta biết được điều đó \mat{$h\approx f$}? (Bài toán cảm ứng \emph{ của Hume })

1) Sử dụng các định lý của lý thuyết học tính toán/thống kê

2) Thử \mat{$h$} trên \defn{bộ thử nghiệm} mới gồm các ví dụ \al
  (sử dụng \emph{cùng phân phối trên không gian mẫu} làm tập huấn luyện)

\defn{Đường cong học tập} = \% đúng trên tập kiểm tra dưới dạng hàm của kích thước tập huấn luyện

\vspace*{-0.2in}

\epsfxsize=0,6\textwidth
\fig{\file{graphs}{restaurant-dtl-curve.ps}}




% Points to make:
% 1) Restaurant data; graph averaged over 20 trials



%%%%%%%%%%%% Slide %%%%%%%%%%%%%%%%%%%%%%%%%%%%%%%%%%%%%%%%%%%%%%%%%%%
\heading{Tiếp tục đo lường hiệu suất }

Đường cong học tập phụ thuộc vào \al
 -- \defn{có thể thực hiện được} (có thể biểu thị hàm mục tiêu) so với \defn{không thể thực hiện được}\nl
     không thể thực hiện được có thể do thiếu thuộc tính \nl
     hoặc lớp giả thuyết bị hạn chế (ví dụ: hàm tuyến tính có ngưỡng)\al
 -- tính biểu đạt dư thừa (ví dụ: vô số thuộc tính không liên quan)

\vspace*{0.2in}

\epsfxsize=0,95\textwidth
\fig{\file{figures}{learning-curves.ps}}



%%%%%%%%%%%% Slide %%%%%%%%%%%%%%%%%%%%%%%%%%%%%%%%%%%%%%%%%%%%%%%%%%%
\heading{Tóm tắt}

Cần học tập cho môi trường chưa rõ, nhà thiết kế lười biếng

Tác nhân học tập = yếu tố hiệu suất + yếu tố học tập

Phương pháp học phụ thuộc vào loại yếu tố hiệu suất, có sẵn\\
phản hồi, loại thành phần cần được cải thiện và cách trình bày của nó

Đối với việc học có giám sát, mục đích là tìm ra một giả thuyết đơn giản\\
điều đó gần như phù hợp với các ví dụ huấn luyện

Học cây quyết định bằng cách sử dụng thông tin đạt được 

Hiệu suất học tập = độ chính xác dự đoán được đo trên bộ kiểm tra

\end{huge}
\end{document}



